% 2.1.7
% Ficheiro incluído por main.tex

\subsection[Utilizadores do AD no Linux: SSSD]{Utilizadores do AD no Linux: SSSD}

\esquema{2-1-7-sssd}{Resolução de utilizadores do AD pelo SSSD}

\passo{Passo 1 - Confirmar que o Linux localiza os utilizadores do AD e ver a configuração do SSSD}

\begin{lstlisting}[style=bashstyle]
id aluno.teste@isec.local
id professor.teste@isec.local
sudo cat /etc/sssd/sssd.conf
\end{lstlisting}

\emph{Verificação esperada:}

\begin{itemize}
\item \texttt{id} mostra o UID e os grupos, ex.: \texttt{groups=...(domain users@isec.local),\allowbreak{}...(grupo-alunos@isec.local)} - os grupos aparecem em \textbf{minúsculas};
\item o \texttt{sssd.conf} (gerado pelo \texttt{realm join}) contém:
\end{itemize}

\begin{center}\small
\begin{tabularx}{\linewidth}{|>{\raggedright\arraybackslash}p{0.32\linewidth}|>{\raggedright\arraybackslash}X|}
\hline
\rowcolor{azul!10}\textbf{Parâmetro} & \textbf{Significado} \\
\hline
\texttt{domains = isec.local} & Domínio que o SSSD serve \\
\hline
\texttt{id\_provider = ad} & Vai buscar utilizadores e grupos ao Active Directory \\
\hline
\texttt{access\_provider = ad} & O AD decide quem pode entrar \\
\hline
\texttt{use\_fully\_qualified\_names = True} & Utilizadores escrevem-se \texttt{utilizador@isec.local} \\
\hline
\texttt{fallback\_homedir = /\allowbreak{}home/\allowbreak{}\%u@\%d} & Homes em \texttt{/\allowbreak{}home/\allowbreak{}utilizador@isec.local} \\
\hline
\end{tabularx}
\end{center}

Se o \texttt{id} der \texttt{no such user}, espera um minuto ou limpa a cache: \texttt{sudo sss\_cache -E} e tenta outra vez.

\problemas{\hyperref[ec-b1-sssd]{Erros Comuns - 2.1 Domínio e SSSD}}
